The appendix is organized into five main sections.
Appendix~\ref{apd:causal_derivations} presents causal derivations and assumptions underlying our analysis. Appendix~\ref{apd:full_results} reports the full experimental results across different model snapshots and causal estimation methods. Appendix~\ref{apd:medication_subgroup_analysis} provides subgroup analyses by medication. Appendix~\ref{apd:dataset_information} describes the study cohort characteristics. Lastly, Appendix~\ref{apd:preprocessing_and_modeling_hyperparameters} details data preprocessing steps and modeling hyperparameters.

\noindent\emph{Each section is further organized as follows:}

In Appendix~\ref{apd:ate-ite}, we derive the relationship between the Average Treatment Effect (ATE) and the Individual Treatment Effect (ITE) under the potential outcomes framework for restricted mean adverse event time.
In Appendix~\ref{adp:from_assoc_to_causal}, we discuss the assumptions required for causal interpretation.
We outline the potential outcomes framework assumptions and the survival analysis assumptions necessary for making causal claims in our setup.
We also briefly highlight challenges with time-varying confounding and suggest directions for improvement.  

Appendix~\ref{apd:causal-meta-learner-algorithms} provides the pseudo-code for the adaptation of causal meta-learners (T-Learner, S-Learner, and Nearest Neighbor Matching Method) to survival analysis.
In Appendix~\ref{apd:matching_derivation} and \ref{apd:csf_ate_derivation}, we illustrate how the treatment effect estimator is derived for the matching meta-learner and Causal Survival Forest respectively. 
Finally, in Appendix~\ref{apd:in_context_of_new_research}, we highlight emerging methodologies that could complement or extend our approach and position our study within a larger research trajectory 

In Appendix~\ref{apd:survival_analysis}, we present the performance of survival analysis models at 9- and 12-month snapshots.  
We summarize model discriminative ability and calibration, highlighting the superiority of the Random Survival Forest.  
In Appendix~\ref{apd:causal_inference}, we report the full causal inference results for 9- and 12-month snapshots.  
We analyze shifts in Average Treatment Effect (ATE) estimates and discuss the role of survival model performance and selection bias. 
The graphical representation of ATE estimates as snapshot time of the prediction is progresses is presented in Appendix~\ref{apd:graphical_ate}.

In Appendix~\ref{apd:unadjusted_survival}, we present Kaplan-Meier survival curves for adherent and non-adherent patients.  
We report restricted mean event time (RMET) estimates as a baseline for later analyses which puts into context that the treatment effects found is not due raw difference of the two cohorts.

In Appendix~\ref{apd:ablation}, we evaluate the effect of removing county-provided risk scores on ATE estimates.  
We analyze how these scores act as proxies for unmeasured confounders across different survival models.  

In Appendix~\ref{apd:ite_medication}, we report Individual Treatment Effect (ITE) distributions for different medications and in Appendix~\ref{apd:ite_injectable_all}, we provide the ITE estimates by medication administration type for all survival model and meta-learner pairs. 

In Appendix~\ref{apd:cohort_information}, we first present a visual overview of the modeling setup and cohort structure (Figure~\ref{fig:data_overview}), illustrating how the first adverse event is defined relative to the prediction snapshot and how monthly adherence is binarized. 
We then describe the study cohort, including demographics, adherence patterns, and prescribing trends.  
We detail demographics in Appendix~\ref{apd:cohort_demographics}, first adverse events timing and types in Appendix~\ref{apd:cohort_first_event},  
non-adherence patterns in Appendix~\ref{apd:cohort_nonadherence}, and antipsychotic medication trends in Appendix~\ref{adp:cohort_medications}.  
In Appendix~\ref{apd:risk_score_evaluations}, we assess the predictive accuracy of county-provided risk scores.  
We present ROC curves and AUC values for predicting specific adverse events within 12 months. 

In Appendix~\ref{apd:preprocessing}, we outline data preprocessing steps.  
We describe cohort filtering, covariate construction, feature normalization, train-test splitting, and trimming strategy for the positivity assumption of our causal analysis.  
In Appendix~\ref{apd:model_hyperparams}, we summarize hyperparameter settings for all survival models.  